% Source-aligned English for printed p.168.
\SourcePage{168}

\LayerTarget{tohoku:proof:3.6.1:p2}
\noindent of subsheaves $F_j$, each of which admits a \textit{finite} family
of generators.
\LayerTarget{tohoku:construction:finite-generator-reduction}
Since $\mathbf C''$ is thick and stable under infinite direct sums, it is also stable
under inductive limits (for an inductive limit of objects in $\mathbf C'$ is
isomorphic to a quotient of their direct sum). Since
$T^p(F)=\varinjlim\limits_j T^p(F_j)$,
to prove $T^p(F)\in \mathbf C''$ it suffices to prove that $T^p(F_j)\in \mathbf C''$ for every
$j$. This reduces us to the case in which $F$ admits a \textit{finite} family
of generators $(f_i)_{1\leqq i\leqq k}$.
\LayerTarget{tohoku:construction:finite-generator-filtration}
Let $F_n$ ($0\leqq n\leqq k$) be the subsheaf of $F$ generated by the $f_i$ with
$1\leqq i\leqq n$. The $F_n$ form a finite increasing sequence of subsheaves of
$F$, each successive quotient $F_n/F_{n-1}$ ($1\leqq n\leqq k$) admitting \textit{one}
generator.
\LayerTarget{tohoku:eq:finite-generator-quotient-reduction}
Induction on the length $k$ of this sequence, together with the fact that
$T^p$ is half-exact, reduces the proof to showing that
$T^p(F_n/F_{n-1})\in \mathbf C''$ for every $n$.
\LayerTarget{tohoku:eq:single-generator-presentation}
We are thereby reduced to the case in which $F$ is generated by a \textit{single}
generator, that is, there is an exact sequence of sheaves
\[
0\to R\to \mathbf Z_U\to F\to0
\]
\LayerTarget{tohoku:eq:single-generator-cohomology}
Here $R$ is a subsheaf of $\mathbf Z_U$, hence also of $\mathbf Z$. We obtain an exact sequence
$T^p(\mathbf Z_U)\to T^p(F)\to T^{p+1}(R)$.
Since $T^p(\mathbf Z_U)\in \mathbf C''$ by hypothesis, it remains to prove that
$T^{p+1}(R)\in \mathbf C''$.
\LayerTarget{tohoku:construction:constant-section-generators}
The subsheaf $R$ of the constant sheaf $\mathbf Z$ is generated by a family
$(f_i)_{i\in I}$ of \textit{constant} sections $n_i$ of $\mathbf Z$ over open subsets $U_i$.
Proceeding as above, we reduce to the case in which this family is finite. We
may of course suppose that $n_i>0$.
\LayerTarget{tohoku:construction:gcd-refinement}
We may further suppose that the $f_i$ have been chosen so that, for every
$x\in\bigcup_i U_i$, among the $n_i$ for which $U_i$ contains $x$ there is a generator of
the subgroup $R(x)$ of $\mathbf Z(x)$, the group $\mathbf Z$ of integers. Indeed, for every
subset $(i_1,\mathinner{.\,.\,.\,.}i_p)$ of $[1,k]$, the set of the first $k$ positive
integers, consider the section $f_{i_1\mathinner{\cdotp\cdotp\cdotp\cdotp}i_p}$ of $R$ over
$U_{i_1}\cap\mathinner{\cdotp\cdotp\cdotp\cdotp}\cap U_{i_p}$ whose value is
$\gcd(n_{i_1},\ldots,n_{i_p})$, and adjoin these sections to the system of
generators.
\LayerTarget{tohoku:construction:R-filtration}
Assume this condition on the $f_i$ and order them so that
$n_1\leqq n_2\leqq\mathinner{.\,.\,.\,.}\leqq n_k$. Let $R_m$ ($0\leqq m\leqq k$) be the subsheaf
of $R$ generated by the sections $f_i$ with $1\leqq i\leqq m$. The $R_m$ form a
finite increasing sequence of subsheaves of $R$.
\LayerTarget{tohoku:eq:R-filtration-reduction}
As above, induction on $k$ and the half-exactness of $T^{p+1}$ show that, to
prove $T^{p+1}(R)\in \mathbf C''$, it suffices to prove
$T^{p+1}(R_m/R_{m-1})\in \mathbf C''$ for $1\leqq m\leqq k$.
\LayerTarget{tohoku:construction:Vm-Ym}
For each $m$, let $V_m$ be the union of the $U_i$ with $1\leqq i\leqq m$, and set
$Y_m=U_m\cap \complement V_{m-1}$.
\LayerTarget{tohoku:eq:R-quotient-ZYm}
We claim that $R_m/R_{m-1}$ is \textit{isomorphic to} $\mathbf Z_{Y_m}$.
\LayerTarget{tohoku:proof:R-quotient-ZYm:p1}
Its restriction to the complement of $V_m$ is plainly zero (since that of
$R_m$ already is), and so is its restriction to $V_{m-1}$. It remains to check
this at each $x\in V_{m-1}\cap U_m$. By hypothesis, the greatest common divisor
$n_{i_0}$ of the $n_i$ corresponding to those $U_i$ that contain $x$ occurs
among them; it therefore divides both $n_m$ and some $n_i$ with $i<m$ (choose
$i<m$ with $x\in U_i$, which exists because $x\in V_{m-1}$). If $i_0<m$, this
shows that $f_m(x)\in R_{m-1}(x)$. If $i_0\geqq m$, then
$n_{i_0}\geqq n_m$, hence $n_{i_0}=n_m$; thus $n_m$, which divides $n_i$ and is
at least $n_i$, equals $n_i$. Again $f_m(x)\in R_{m-1}(x)$. In both cases
$R_{m-1}(x)=R_m(x)$. Hence the restriction of $R_m/R_{m-1}$ to the union of
$\complement V_m$ and $V_{m-1}$, that is, to $\complement Y_m$, is

\TohokuContinuation{169}
